\documentclass[10pt,a4paper]{amsart}
\usepackage[margin=19mm]{geometry}
\usepackage{amsmath,amssymb,mathtools}
\usepackage[scr=rsfs,cal=euler]{mathalfa}
\usepackage{xfrac}
\usepackage[unicode,hidelinks,bookmarks=false]{hyperref}
\newcommand{\ie}{i.e.\ }
\newcommand{\cf}{cf.\ }
\newcommand{\resp}{respectively\ }
\newcommand{\Subsection}{\subsection}
\providecommand{\nobreakdash}{\nobreak\hbox{-}}
\setlength{\emergencystretch}{2em}
\multlinegap0pt
\usepackage{bm}
\usepackage{bbm}
\usepackage{dsfont}
\usepackage{xspace}
\usepackage{cancel}
\usepackage{mathtools}
\usepackage{amsthm}
\makeatletter
\@ifundefined{thm}{\newtheorem{thm}{Thm}}{}
\@ifundefined{lem}{\newtheorem{lem}[thm]{Lemma}}{}
\makeatother
\DeclareMathOperator{\nai}{INT}
\newcommand{\metdis}{\mathcal{D}}
\newcommand{\met}[1]{\mathrm{Met}(#1)}
\newcommand{\nn}{\mathbb{N}}
\newcommand{\yodset}[1]{B(#1)}
\newcommand{\yoeset}[1]{\mathrm{E}_{#1}}
\newcommand{\yokset}[1]{K(#1)}
\newcommand{\yolset}[1]{\mathrm{L}_{#1}}
\newcommand{\yopset}{\mathrm{PM}}
\title[An interpolation of metrics and spaces of metrics]{An interpolation of metrics and spaces of metrics}
\author{Yoshito Ishiki}
\date{Source: arXiv:2003.13227v3 (CC BY 4.0)}
\begin{document}
\maketitle

\noindent\emph{Source and licence.} An interpolation of metrics and spaces of metrics. Version arXiv:2003.13227v3 (CC BY 4.0), distributed under the Creative Commons Attribution 4.0 International licence, as recorded in the arXiv licence metadata for this exact version and shown on the abstract page https://arxiv.org/abs/2003.13227v3. Full rights notice: licenses/arxiv-2003.13227-CC-BY-4.0.txt.\par\smallskip

\noindent\textit{Editorial scope.} Counted unit: the Lem whose source label is \texttt{lem:Baire} in the article (source0.tex lines 3214--3222, source label \texttt{lem:Baire}), delivered together with the article's own complete original proof of it (source0.tex lines 3223--3579, one \texttt{proof} environment of 357 lines). No mathematics is written, repaired or re-derived: statement and proof are the article's own text line for line. Required context retained uncounted: lem:gdelta (lines 1118--1122). Every definition, display and label of those lines is retained unchanged. Omitted from this package and not redistributed: 1--1117, 1123--3213, 3580--3692. Nothing inside a retained excerpt is omitted or altered: every retained body is delivered exactly as the article's lines are written, so no omission receipt of any kind accompanies this package. Citations: the retained excerpts carry no citation command at all, so no bibliography entry of the article is transcribed and no reference is redistributed. Reference adaptation: none was needed and none was made; every reference inside the delivered unit resolves inside it, so no unresolved reference or citation is produced. Printed slips kept literal: no printed word, symbol, hypothesis or number of the counted statement or of its proof was altered. Layout and class: the article's own class is \texttt{amsart} with packages including amssymb, amsthm, amsmath, eucal, mathrsfs, comment, enumitem; the wrapper is the lane's pinned \texttt{amsart} editorial wrapper at 10pt on A4, which restates the theorem-like environments the retained text needs and adds the article's own macro definitions that the retained text needs (\texttt{\textbackslash met}, \texttt{\textbackslash metdis}, \texttt{\textbackslash nai}, \texttt{\textbackslash nn}, \texttt{\textbackslash yodset}, \texttt{\textbackslash yoeset}, \texttt{\textbackslash yokset}, \texttt{\textbackslash yolset}, \texttt{\textbackslash yopset}), with extra wrapper packages (bm, bbm, dsfont, xspace, cancel, mathtools). The delivered numbering is the unit's own. Assets: no retained span contains a figure environment, a graphics-inclusion command, a PDF-page inclusion, a table or a TikZ picture; no image file is copied into this package, the package declares no asset dependency and no image file is redistributed with it. Licence: version arXiv:2003.13227v3 of this article is distributed under the Creative Commons Attribution 4.0 International licence, as recorded in the arXiv licence metadata for this exact version and shown on the abstract page https://arxiv.org/abs/2003.13227v3. A full rights notice is delivered at licenses/arxiv-2003.13227-CC-BY-4.0.txt. Characters: every retained excerpt is pure ASCII, so the delivered engine, XeLaTeX with its default Latin Modern text font, reports no missing character.
\begin{lem}\label{lem:gdelta}
Every 
$G_{\delta}$ 
subset of a completely metrizable space is a Baire space. 
\end{lem}

\begin{lem}\label{lem:Baire}
For 
every  second-countable  locally compact metrizable space 
$X$, 
the space 
$\met{X}$
 is  completely metrizable. 
 In particular,  it is Baire. 
\end{lem}

\begin{proof}
Let 
$C(X^{2})$ 
be the set of all real-valued continuous functions on $X^{2}$.
We define a metric 
$\mathcal{E}$ 
on
 $C(X^{2})$ 
by 
\begin{equation*}\label{eq:E}
\mathcal{E}(f, g)=\min\left\{1,\sup_{(x, y)\in X^{2}}|f(x, y)-g(x, y)|\right\}. 
\end{equation*}
Note that
the metric 
$\mathcal{E}|_{\met{X}^{2}}$ 
on 
$\met{X}$  
generates the same topology as  
$\metdis_{X}$ 
on 
$\met{X}$, 
which coincides  with the uniform topology.  




Note that  the space 
$(C(X^{2}), \mathcal{E})$ 
is completely metrizable. 
By Lemma 
\ref{lem:gdelta}, 
to prove the lemma, 
it suffices to show that 
$\met{X}$ 
is 
$G_{\delta}$ 
in  
$(C(X^{2}), \mathcal{E})$.
We denote by
 $\yopset$ 
 the set of all 
$f\in C(X^{2})$ 
such that 
%%%%%
\begin{enumerate}
	
	%%%%%
	\item for every 
	$x\in X$ 
	we have 
	$f(x, x)=0$;
	
	%%%%%
	\item 
	for all
	 $x, y\in X$ 
	 we have 
	 $f(x, y)=f(y, x)$
	 and 
	 $f(x, y)\ge 0$;
	 
	 %%%%%%
	\item for every triple 
	$x, y, z\in X$
	 we have  $
	 f(x, y)\le f(x, z)+f(z, y)$. 
\end{enumerate}
%%%%%%%%%%%%%
%%%%%%%%%%%%
%%%%%%%%%%%%%
Namely, 
$\yopset$ 
is the set of all 
continuous pseudo-metrics on 
$X$. 
Note that 
$\yopset$ 
is 
a closed subset in  the metric space 
$(C(X^{2}), \mathcal{E})$. 
Since all closed subsets 
of a metric space are 
$G_{\delta}$ 
in the whole space, 
the set 
$\yopset$ 
is 
$G_{\delta}$ 
in 
$(C(X^{2}), \mathcal{E})$. 


Using the assumption that 
$X$ 
is second-countable and  locally compact, 
now  we take a sequence 
$\{\yodset{n}\}_{n\in \nn}$ 
of compact subsets of 
$X^{2}$ 
with 
$\bigcup_{n\in \nn}\yodset{n}=X^{2}\setminus \Delta_{X}$, 
where 
$\Delta_{X}$ 
is the diagonal set of 
$X^{2}$, 
and take a sequence 
$\{\yokset{n}\}_{n\in \nn}$ 
of compact subsets of 
$X$ 
with 
$\yokset{n}\subset \nai(\yokset{n+1})$ 
and 
$\bigcup_{n\in \nn} \yokset{n}=X$, 
where 
``$\nai$''
means the interior. 


%%%
For every 
$n\in \nn$,  
let
 $\yolset{n}$ 
be the set of all 
$f\in C(X^{2})$ 
for which 
there exists
$c\in (0, \infty)$ 
such that 
for each  
$s>n$ 
we have 
\[
\inf\{\, f(x, y)\mid x\in \yokset{n}, \  y\in X\setminus \yokset{s}\, \} >c. 
\]

For each 
$n\in \nn$,  
let 
$\yoeset{n}$ 
be 
the set of all 
$f\in C(X^{2})$ 
such that
for each 
$(x, y)\in \yodset{n}$ 
we have 
$0<f(x, y)$.  
Note that each 
$\yolset{n}$ 
and each 
$\yoeset{n}$ 
are open subsets 
of 
$(C(X^{2}), \mathcal{E})$. 

We next prove that
\[
	\met{X}=
	\yopset
	\cap 
	\left(
	\bigcap_{n\in \nn} \yolset{n}
	\right)
	\cap  
	\left(
	\bigcap_{n\in \nn}\yoeset{n}
	\right). 
\]
To prove
$\met{X}\subset
	\yopset
	\cap 
	\left(
	\bigcap_{n\in \nn} \yolset{n}
	\right)
	\cap  
	\left(
	\bigcap_{n\in \nn}\yoeset{n}
	\right)$, 
take 
$d\in \met{X}$. 
Since 
$d$
 is a metric on 
 $X$, 
we have
 $d\in \yopset$ 
and  
$d\in \bigcap_{n\in \nn}\yoeset{n}$. 
To show 
$d\in \bigcap_{n\in \nn}\yolset{n}$,  
take an arbitrary number 
$n\in \nn$. 
Since 
$\yokset{n}\subset \nai(\yokset{n+1})$ 
and 
since 
$d\in \met{X}$, 
for every 
$s> n$ 
we have 
\begin{align*}
	0&<d(\yokset{n}, X\setminus \nai(\yokset{n+1}))
	\le
	d(\yokset{n}, X\setminus \yokset{n+1})
	\\
	&\le
	d(\yokset{n}, X\setminus \yokset{s}).
\end{align*}
Thus  
$d\in \bigcap_{n\in \nn}\yolset{n}$. 
Hence we obtain
\[
	\met{X}
	\subset 
	\yopset
	\cap 
	\left(
	\bigcap_{n\in \nn} \yolset{n}
	\right)
	\cap
	\left(
	\bigcap_{n\in \nn}\yoeset{n}
	\right). 
\]
Next we  take 
$d\in \yopset\cap \left(\bigcap_{n\in \nn} \yolset{n}\right)\cap  \left(\bigcap_{n\in \nn}\yoeset{n}\right)$. 
Since 
$d\in \yopset\cap \left( \bigcap_{n\in \nn}\yoeset{n}\right)$, 
the function 
$d$ 
is  continuous on 
$X^{2}$ 
and it is  a metric on 
$X$. 
Fix 
$e\in \met{X}$. 
We show that  
the metric  
$d$ 
is topologically equivalent to
 $e$. 
Since 
$d$ 
is continuous on 
 $X^{2}$, 
the metric
 $d$ 
 generates a weaker topology than that of
  $(X, e)$. 
Namely, 
if 
$\lim_{n\to \infty}x_{n}= a$ 
in 
$(X, e)$, 
then 
$\lim_{n\to \infty}x_{n}=a$ 
in 
$(X, d)$. 
Assume next  that 
$\lim_{n\to \infty}x_{n}= a$
 in 
 $(X, d)$. 
Our purpose is 
to prove 
that 
$\lim_{n\to \infty}x_{n}\to a$
in 
$(X, e)$. 
Since 
$\{\yokset{n}\}_{n\in \nn}$ 
is a covering of 
$X$, 
there exists 
$s\in \nn$
such that 
$a\in \yokset{s}$. 
For the sake of contradiction, 
suppose that 
there exist infinitely many 
$i\in \nn$ 
with
 $(X\setminus \yokset{i})\cap \{\, x_{n}\mid {n\in \nn}\, \}\neq \emptyset$. 
Then we have
\[
	\liminf_{i\to \infty}d(\yokset{s}, X\setminus \yokset{i})
	\le 
	\lim_{j\to \infty}d(a, x_{j})
	=
	0.
\] 
This contradicts the fact that 
$d\in \bigcap_{n\in \nn}\yolset{n}$. 
Hence there exists 
a sufficiently large number 
$m\in \nn$ 
such that  
the sequence 
$\{x_{n}\}_{n\in \nn}$ 
is contained in 
$\yokset{m}$. 
 We may assume that 
 $s\le m$, which implies that 
 $a\in \yokset{m}$. 
By the compactness of 
$\yokset{m}$ and by 
the fact that 
$(X, d)$ 
is Hausdorff, 
the map 
$1_{\yokset{m}}\colon (\yokset{m}, e)\to (\yokset{m}, d)$
becomes a homeomorphism. 
Thus 
the restricted metrics 
$e|_{\yokset{m}^{2}}$ and 
$d|_{\yokset{m}^{2}}$ 
generate the same topology on 
$\yokset{m}$. 
Since  
$\lim_{n\to\infty}x_{n}=a$ in $(X, d)$
and since 
$\{x_{n}\}_{n\in \nn}$
 is 
 contained in 
$\yokset{m}$, 
we see that 
$\lim_{n\to\infty}x_{n}=a$
in 
$(X, e)$. 

As a result, 
the metric 
$d$
 generates the same topology as 
 $e$,  
and hence 
\[
	\met{X}
	\supset 
	\yopset
	\cap 
	\left(
	\bigcap_{n\in \nn} \yolset{n}
	\right)
	\cap  
	\left(
	\bigcap_{n\in \nn}\yoeset{n}
	\right). 
\]
Therefore we conclude that 
$\met{X}$ 
is a 
$G_{\delta}$ 
subset of 
$(C(X^{2}), \mathcal{E})$. 
\end{proof}

\end{document}
